\section{Complete multiplicative seeds at the golden base}
\label{golden:sec:complete-seeds}

The complementary-power classification in the current manuscript
\cite{ProveIt} concerns relations of the form $r^a+r^b=1$.
Golden ladders also use products of many factors $1-r^n$.
We now give a completeness theorem for all such finite multiplicative
relations at the golden base, without imposing an exponent bound in
advance.  The result determines an integral relation lattice, explains
the eight arguments in the manuscript's higher golden formulas, and
classifies which integer powers of the base retain any such seeds.

The four exceptional indices $1,6,10,12$, when the base is the square
of the golden reciprocal, were already recorded by Bailey and
Broadhurst \cite[\S2, preceding (22)]{BaileyBroadhurst1999}.
We do not claim discovery of that list.  The contribution here is its
unbounded completeness proof, an explicit integral lattice for all
positive powers, and its consequences for power restrictions and
exterior symbols.  The arithmetic input is Flatters' later theorem
on primitive divisors \cite{Flatters2009}.

\subsection{An effective bound for every real quadratic unit}
\label{golden:subsec:quadratic-cutoff}

Let $K$ be a real quadratic field, let $\sigma$ denote its nontrivial
automorphism, and let $r\in\mathcal O_K^\times$ satisfy $0<r<1$.
Put $V_n(r)=1-r^n$ for $n\ge1$.
A prime ideal $\mathfrak p$ is \emph{primitive for $V_n(r)$} if it
divides $V_n(r)$ and none of $V_1(r),\ldots,V_{n-1}(r)$.
Because $r$ is a unit, this says exactly that its residue in the finite
field $\mathcal O_K/\mathfrak p$ has multiplicative order $n$.

We use the following part of Flatters' theorem: if $u$ is a positive
real quadratic unit of norm $+1$, then
$N_{K/\mathbb Q}(u^m-1)$ has a primitive rational prime divisor for
every $m>12$ \cite[Theorem~1.4]{Flatters2009}.
Here a rational prime is primitive if it divides this norm and no
earlier nonzero norm.  The following passage to prime ideals is needed
when the original unit has norm $-1$.

\begin{theorem}[Uniform quadratic cutoff]
\label{golden:thm:quadratic-cutoff}
If $N_{K/\mathbb Q}(r)=+1$, then $1-r^n$ has a primitive prime-ideal
divisor for every $n>12$.  If $N_{K/\mathbb Q}(r)=-1$, it has one
for every $n>24$.

Consequently, in any finite relation
\begin{equation}
 \prod_{n\ge1}(1-r^n)^{c_n}=r^b,
 \qquad c_n,b\in\mathbb Z,
 \label{golden:eq:general-seed}
\end{equation}
all coefficients above $12$ vanish in the norm $+1$ case, and all
coefficients above $24$ vanish in the norm $-1$ case.
The same support conclusion holds for rational exponents.
\end{theorem}

\begin{proof}
For norm $+1$, apply Flatters' theorem to $u=r^{-1}>1$.
At index $n>12$, choose a primitive rational divisor $p$ of
$N_{K/\mathbb Q}(u^n-1)$ and a prime ideal $\mathfrak p$ above $p$
dividing $u^n-1$.  The residue order of $u$ is exactly $n$: an order
$j<n$ would make $p$ divide the earlier norm
$N_{K/\mathbb Q}(u^j-1)$.  Inversion preserves order, so $r$ also
has order $n$ at this ideal.

For norm $-1$, put $q=r^2$ and apply the norm $+1$ theorem to
$u=r^{-2}$.  For each $m>12$, we obtain a primitive rational divisor
of $N_{K/\mathbb Q}(q^m-1)$, because
$u^j-1=-q^{-j}(q^j-1)$ differs by a unit.  At a suitable ideal
$\mathfrak p$ above this prime, the residue of $q$ has exact order
$m$.

The residue characteristic is odd.  Indeed, the residue field of an
ideal above $2$ in a quadratic field has size $2$ or $4$; the residue
of any unit therefore has order $1$ or $3$.  Thus $2$ divides some
norm $N_{K/\mathbb Q}(q^j-1)$ with $j\le3$, and cannot be
primitive at $m>12$.

Let $a$ be the residue of $r$ at $\mathfrak p$.  Since $a^2$ has
order $m$, the order of $a$ is either $m$ or $2m$, and it must be
$2m$ when $m$ is even.  The conjugation formula
\begin{equation}
 \sigma(r)=-r^{-1}
 \label{golden:eq:normminus-conjugation}
\end{equation}
identifies the residue of $r$ at the conjugate ideal, under the
induced field isomorphism, with $-a^{-1}$.  When $m$ is odd, these
two residues have orders $m$ and $2m$ in some order.  To check this
directly, if $a$ has odd order $m$, then $-a^{-1}$ has order $2m$.
If $a$ has order $2m$, then $a^m=-1$ in odd characteristic, so
$(-a^{-1})^m=1$; its square still has order $m$, forcing its order
to be $m$.
In this odd-$m$ case the ideals must be distinct: if they coincided,
the induced residue-field automorphism would preserve order.  Thus
no splitting assumption is needed.

For an even target $n>24$, take $m=n/2>12$.  If $m$ is even, the
chosen ideal gives residue order $n$.  If $m$ is odd, the chosen
ideal or its conjugate gives order $2m=n$.  For an odd target
$n>24$, take $m=n$ and choose the ideal giving order $m=n$.
This proves the prime-ideal statement in every case.

Finally, suppose \eqref{golden:eq:general-seed} has a nonzero
coefficient above the relevant cutoff, and choose its largest
supported index $N$.  At a primitive ideal for $V_N(r)$, all lower
factors and the unit $r$ have valuation zero.  Taking valuations gives
$c_N v_{\mathfrak p}(V_N(r))=0$, which contradicts both $c_N\ne0$
and $v_{\mathfrak p}(V_N(r))>0$.  Clearing denominators proves the
rational-exponent version.
\end{proof}

This theorem bounds multiplicative seeds.  It does not place an
unconditional bound on the weight of all polylogarithmic identities.
For a fixed quadratic base, it reduces seed determination to finitely
many exact ideal valuations and a unit calculation.

\subsection{The full integral lattice at the golden reciprocal}
\label{golden:subsec:integral-lattice}

Set
\[
 \rho=\frac{\sqrt5-1}{2},\qquad V_n=1-\rho^n,
 \qquad K=\mathbb Q(\sqrt5),\qquad\mathcal O_K=\mathbb Z[\rho].
\]
All $V_n$ are positive, so the following equations have no real
logarithm or sign ambiguity.

\begin{theorem}[Complete golden seed lattice]
\label{golden:thm:integral-lattice}
The group of finite integer relations
\begin{equation}
 \prod_{n\ge1}V_n^{c_n}=\rho^b,
 \qquad (c_n)\in\mathbb Z^{(\mathbb Z_{>0})},\quad b\in\mathbb Z,
 \label{golden:eq:rho-seed}
\end{equation}
is free of rank six.  An integral basis is
\begin{equation}
\begin{aligned}
 V_1&=\rho^2,&
 V_2&=\rho,\\
 V_6&=\rho^{-1}V_3^2,&
 V_{12}&=\rho^{-2}V_4V_3^3,\\
 V_{20}&=\rho^{-1}V_4^3V_{10},&
 V_{24}V_4&=\rho^{-2}V_3^4V_8^2.
\end{aligned}
\label{golden:eq:six-generators}
\end{equation}
Equivalently, $c$ occurs in \eqref{golden:eq:rho-seed} precisely when
its support is contained in
\begin{equation}
 S=\{1,2,3,4,6,8,10,12,20,24\}
 \label{golden:eq:full-support}
\end{equation}
and
\begin{equation}
\begin{aligned}
 c_3+2c_6+3c_{12}+4c_{24}&=0,\\
 c_4+c_{12}+3c_{20}-c_{24}&=0,\\
 c_8+2c_{24}&=0,\\
 c_{10}+c_{20}&=0.
\end{aligned}
\label{golden:eq:four-constraints}
\end{equation}
The uniquely determined unit exponent is
\begin{equation}
 b=2c_1+c_2-c_6-2c_{12}-c_{20}-2c_{24}.
 \label{golden:eq:unit-exponent}
\end{equation}
\end{theorem}

We first record the exact finite information needed in the proof.
Reducing powers using $\rho^2=1-\rho$ gives
\begin{center}
\small
\begin{tabular}{c|c@{\qquad}c|c}
$n$ & $V_n$ & $n$ & $V_n$\\ \hline
$1$ & $\rho^2$ & $2$ & $\rho$\\
$3$ & $2\rho^2$ & $4$ & $\sqrt5\rho^2$\\
$6$ & $4\rho^3$ & $8$ & $3\sqrt5\rho^4$\\
$10$ & $11\rho^5$ & $12$ & $8\sqrt5\rho^6$\\
$20$ & $55\sqrt5\rho^{10}$ & $24$ & $144\sqrt5\rho^{12}$
\end{tabular}
\end{center}
Substitution proves every equation in
\eqref{golden:eq:six-generators}.  In particular, the indices
\[
 E=\{1,2,6,12,20,24\}
\]
introduce no prime ideal absent from lower indices.

Every other index through $24$ has the following primitive witness.
A scalar $a$ denotes the ideal $(p,\rho-a)$ with residue field
$\mathbb F_p$.  The entry $\theta$ denotes the ideal $(p)$ with
residue field
$\mathbb F_p[\theta]/(\theta^2+\theta-1)$; the defining polynomial
is irreducible for the three primes $p=2,3,7$ used below.
All powers in the last column are computed in the indicated residue
field.

\begin{center}
\small
\begin{tabular}{r|r|r|l}
$n$ & $p$ & residue of $\rho$ & proper-power checks\\ \hline
$3$ & $2$ & $\theta$ & $\theta^1=\theta$\\
$4$ & $5$ & $2$ & $2^2=4$\\
$5$ & $11$ & $3$ & $3^1=3$\\
$7$ & $29$ & $23$ & $23^1=23$\\
$8$ & $3$ & $\theta$ & $\theta^4=2$\\
$9$ & $19$ & $4$ & $4^3=7$\\
$10$ & $11$ & $7$ & $7^5=10,\quad 7^2=5$\\
$11$ & $199$ & $61$ & $61^1=61$\\
$13$ & $521$ & $421$ & $421^1=421$\\
$14$ & $29$ & $5$ & $5^7=28,\quad5^2=25$\\
$15$ & $31$ & $18$ & $18^5=25,\quad18^3=4$\\
$16$ & $7$ & $\theta$ & $\theta^8=6$\\
$17$ & $3571$ & $1103$ & $1103^1=1103$\\
$18$ & $19$ & $14$ & $14^9=18,\quad14^6=7$\\
$19$ & $9349$ & $7563$ & $7563^1=7563$\\
$21$ & $211$ & $178$ & $178^7=196,\quad178^3=144$\\
$22$ & $199$ & $137$ & $137^{11}=198,\quad137^2=63$\\
$23$ & $139$ & $63$ & $63^1=63$
\end{tabular}
\end{center}

For each scalar row, direct arithmetic verifies
$a^2+a-1=0$ and $a^n=1$ modulo $p$; for each quadratic row it
verifies $\theta^n=1$.  The last column lists the values at $n/\ell$
for every prime $\ell\mid n$, and none equals $1$.
Thus the order is exactly $n$.  The listed rational primes are
verified by trial division, and irreducibility of the three quadratic
polynomials is checked by testing their residues.  The accompanying
script replays all of these finite checks in exact arithmetic.

\begin{proof}[Proof of Theorem~\ref{golden:thm:integral-lattice}]
Theorem~\ref{golden:thm:quadratic-cutoff} shows that every relation
\eqref{golden:eq:rho-seed} is supported on indices at most $24$.
Subtract integer multiples of the six rows in
\eqref{golden:eq:six-generators} to set the coefficients at
$1,2,6,12,20,24$ to zero.  Each row has coefficient $+1$ in its
own distinct pivot column and zero in every other pivot column, so
this elimination is integral and leaves all eliminated coefficients
zero.

If the residual relation is nonzero, let $N$ be its largest supported
index.  Then $N\notin E$ and the table supplies an ideal primitive
for $V_N$.  Every lower factor and $\rho$ has valuation zero there,
whereas $V_N$ has positive valuation.  Applying this valuation again
contradicts the nonzero coefficient at $N$.  Hence the six rows span
the entire relation group over $\mathbb Z$.

Their distinct pivot columns also prove independence.  Reading their
other columns gives \eqref{golden:eq:four-constraints}, and reading
their unit exponents gives \eqref{golden:eq:unit-exponent}.  The unit
exponent is unique because $\rho$ is not a root of unity.
\end{proof}

\begin{remark}[Why the ideal witnesses matter]
The norms of $V_5$ and $V_{10}$ are $11$ and $-121$.
One ideal above $11$ has residue $\rho=3$ of order $5$; its conjugate
has residue $\rho=7$ of order $10$.  Thus $V_{10}$ has a primitive
prime-ideal divisor even though its norm has no primitive rational
prime divisor.  A proof using only new rational prime factors of the
norms would miss this constraint.
\end{remark}

\subsection{Four even-power seeds and an integral parity condition}
\label{golden:subsec:even-seeds}

Put $q=\rho^2$ and $U_j=1-q^j=V_{2j}$.

\begin{corollary}[The complete even-power lattice]
\label{golden:cor:even-seeds}
All finite integral relations
$\prod_{j\ge1}U_j^{c_j}=\rho^b$ are generated freely by
\begin{equation}
\begin{aligned}
 U_1&=\rho,\\
 U_6^2&=\rho^{-1}U_2^2U_3^3,\\
 U_{10}&=\rho^{-1}U_2^3U_5,\\
 U_{12}U_2&=U_3^2U_4^2.
\end{aligned}
\label{golden:eq:four-even-generators}
\end{equation}
Their support is contained in
\begin{equation}
 \{1,2,3,4,5,6,10,12\}
 \quad\text{in the $q$ index}.
 \label{golden:eq:even-support}
\end{equation}
If the four rows are used with multiplicities $(A,B,C,D)$,
the resulting exponent of $\rho$ is $A-B-C$.
The result is an integer power of $q$ precisely when $A-B-C$ is even.
\end{corollary}

\begin{proof}
Write the multiplicities of the rows in
\eqref{golden:eq:six-generators} as
$a_1,a_2,a_6,a_{12},a_{20},a_{24}$.
Their odd-index coefficients vanish exactly when
\[
 a_1=0,\qquad 2a_6+3a_{12}+4a_{24}=0.
\]
Thus $a_{12}$ is even and
$a_6=-3(a_{12}/2)-2a_{24}$.
The four free integer parameters
$a_2,a_{12}/2,a_{20},a_{24}$ give exactly
\eqref{golden:eq:four-even-generators}.  The unit exponent follows
from the four displayed identities, proving the parity assertion.
\end{proof}

The support \eqref{golden:eq:even-support} is precisely the
eight-argument set
\begin{equation}
 \mathcal A=\{\rho^2,\rho^4,\rho^6,\rho^8,
                 \rho^{10},\rho^{12},\rho^{20},\rho^{24}\}
 \label{golden:eq:manuscript-alphabet}
\end{equation}
used by the manuscript's conjectured higher-weight golden ladders.
The theorem now establishes that no arbitrarily large even exponent
can introduce another multiplicative seed.  At the level of rational
logarithmic relations, the integral parity condition disappears because
$\mathbb Q\log q=\mathbb Q\log\rho$.

\subsection{Which powers of the base retain any seeds?}
\label{golden:subsec:power-restriction}

For a positive integer $h$, define the rational relation space
\begin{equation}
 \mathcal R_h=
 \left\{(c_n)\in\mathbb Q^{(\mathbb Z_{>0})}:
  \sum_{n\ge1}c_n\log(1-\rho^{hn})
      \in\mathbb Q\log\rho\right\}.
 \label{golden:eq:restricted-space}
\end{equation}
One may replace the line $\mathbb Q\log\rho$ by
$\mathbb Q\log(\rho^h)$ without changing this space.

\begin{theorem}[Sharp power-restriction ranks]
\label{golden:thm:power-ranks}
For every positive integer $h$,
\begin{equation}
 \dim_{\mathbb Q}\mathcal R_h=
 \begin{cases}
 6,&h=1,\\
 4,&h=2,\\
 1,&h=3,4,\\
 0,&h\ge5.
 \end{cases}
 \label{golden:eq:power-ranks}
\end{equation}
For $h=3$, the unique rational direction is
$V_6/V_3^2=\rho^{-1}$.
For $h=4$, it is
\begin{equation}
 \frac{V_{12}^4V_8^6}{V_4^7V_{24}^3}=\rho^{-2}.
 \label{golden:eq:power-four-direction}
\end{equation}
The primitive integral relations whose right sides are integer powers
of their own base $r=\rho^h$ are
\begin{align}
 h=3:\quad &(1-r^2)^3=r^{-1}(1-r)^6,
 \label{golden:eq:power-three-integral}\\
 h=4:\quad &(1-r^3)^8(1-r^2)^{12}
       =r^{-1}(1-r)^{14}(1-r^6)^6.
 \label{golden:eq:power-four-integral}
\end{align}
\end{theorem}

\begin{proof}
Clear denominators in \eqref{golden:eq:restricted-space} and
exponentiate.  The resulting relation has the form
\eqref{golden:eq:rho-seed} and is supported only on multiples of $h$.
Theorem~\ref{golden:thm:integral-lattice} supplies the complete
support and constraints.

The cases $h=1,2$ follow immediately from that theorem and
Corollary~\ref{golden:cor:even-seeds}.
For $h=3$, the permitted support is $\{3,6,12,24\}$.
The equation at index $8$ in \eqref{golden:eq:four-constraints}
forces $c_{24}=0$; the equation at index $4$ then forces
$c_{12}=0$.  The remaining freedom is $c_3=-2c_6$.

For $h=4$, the permitted support is $\{4,8,12,20,24\}$.
The equation at index $10$ forces $c_{20}=0$; that at index $3$
gives $3c_{12}+4c_{24}=0$.
Thus the integer solutions are
\[
 c_{12}=4t,\qquad c_{24}=-3t,\qquad
 c_4=-7t,\qquad c_8=6t,
\]
which give \eqref{golden:eq:power-four-direction}.
Their unit exponent is $-2t$, so it is an integer power of
$\rho^4$ precisely when $t$ is even.  For $h=3$, the unit
exponent is $-c_6$, so an integer power of $\rho^3$ requires
$3\mid c_6$.  These conditions give the two primitive integral
relations displayed in the theorem.

For $h\ge5$, support is empty unless $h$ divides a member of
$S$ in \eqref{golden:eq:full-support}.  The possibilities are
$h\in\{5,6,8,10,12,20,24\}$.  It suffices to exclude $5,6,8$,
as their multiples impose further restrictions.
For $h=5$, only indices $10,20$ are permitted; the index-$4$
equation forces $c_{20}=0$, and the index-$10$ equation then
forces $c_{10}=0$.
For $h=6$, only $6,12,24$ are permitted; the equations at indices
$8,4,3$ eliminate them in that order.
For $h=8$, only $8,24$ are permitted; the equation at index $4$
eliminates $24$, then that at index $8$ eliminates $8$.
This proves all cases of \eqref{golden:eq:power-ranks}.
\end{proof}

The number field remains $\mathbb Q(\sqrt5)$ under these power
substitutions, but the permitted positive-power seed space can shrink
from dimension six to zero.  Thus the field alone does not determine
the seed count for a specified base and argument family.

\subsection{All-weight kernels of the exterior-symbol map}
\label{golden:subsec:exterior-symbol}

Put
\[
 W=K^\times\otimes_{\mathbb Z}\mathbb Q,
 \qquad\langle x\rangle=x\otimes1\in W,
 \qquad v=\langle\rho\rangle,
\]
and use additive notation for $W$.
For $w\ge2$, define the rationalized classical exterior-symbol map
\begin{equation}
 \beta_w([x])=
 \langle x\rangle^{\otimes(w-2)}\otimes
 \bigl(\langle x\rangle\wedge\langle1-x\rangle\bigr)
 \quad (x\ne0,1),
 \label{golden:eq:beta-definition}
\end{equation}
with target $W^{\otimes(w-2)}\otimes\bigwedge^2W$.
Here $[x]$ denotes a generator in a free rational vector space.
We make no identification of its quotient with numerical periods.

\begin{theorem}[Complete positive-power exterior-symbol kernels]
\label{golden:thm:symbol-kernel}
On the free span of $[\rho^n]$, $n\ge1$, the kernel of $\beta_w$
has dimension six at every weight $w\ge2$.
An explicit basis consists of
\begin{equation}
 \sum_n\frac{c_n}{n^{w-1}}[\rho^n],
 \label{golden:eq:weighted-seed}
\end{equation}
where $c$ runs through the six coefficient vectors in
\eqref{golden:eq:six-generators}.
Every kernel element is supported on
\eqref{golden:eq:full-support}.
On even powers the kernel has dimension four and is supported on
\eqref{golden:eq:manuscript-alphabet}.
On the span of $[\rho^{hn}]$, its dimensions are
\eqref{golden:eq:power-ranks}.
\end{theorem}

\begin{proof}
The vector $v$ is nonzero because $\rho$ is not torsion.
Multilinearity gives, for every finitely supported rational vector $a$,
\begin{equation}
 \beta_w\!\left(\sum_n a_n[\rho^n]\right)
 =v^{\otimes(w-2)}\otimes
 \left(v\wedge\sum_n a_n n^{w-1}\langle V_n\rangle\right).
 \label{golden:eq:beta-factorization}
\end{equation}
Tensoring over a field with a fixed nonzero pure tensor is injective,
and $v\wedge z=0$ if and only if $z\in\mathbb Qv$.
Consequently \eqref{golden:eq:beta-factorization} vanishes exactly
when
\begin{equation}
 \sum_n a_n n^{w-1}\langle V_n\rangle\in\mathbb Qv.
 \label{golden:eq:weighted-line-test}
\end{equation}

After clearing denominators, this says that a positive real algebraic
number $\prod_nV_n^{c_n}\rho^{-b}$ is torsion.  A positive real
root of unity is $1$, so the condition is equivalent to the actual
multiplicative relation \eqref{golden:eq:rho-seed}, with the
coefficients proportional to $a_n n^{w-1}$.
The diagonal operation $a_n\mapsto n^{w-1}a_n$ is invertible and
preserves support.  Theorems~\ref{golden:thm:integral-lattice}
and \ref{golden:thm:power-ranks} therefore give the basis, support,
and every stated dimension.
\end{proof}

\paragraph{Scope of the obstruction.}
This theorem computes the kernel of the explicitly specified algebraic
map.  A nonzero exterior symbol rules out a certificate requiring
vanishing in this map.  Vanishing in the number-field symbol alone does
not supply a rational-function deformation, lower-weight compatibility,
or the regulator information needed to prove a higher polylogarithm
identity.  In particular, dimension four does not assert that four
numerical ladders exist at every weight, and it does not prove that
none exists above weight nine.  Numerical period independence remains
a separate arithmetic question.

The constructive consequence is that searches beginning with this
symbol can use \eqref{golden:eq:manuscript-alphabet} as a complete
even-power alphabet and work in four weighted seed coordinates.
Increasing the positive even exponents beyond $24$ cannot enlarge that
particular kernel.

\subsection{Exact replay and next proof problems}
\label{golden:subsec:replay}

The file \texttt{code/golden\_seed\_certificate.py} uses only
Python's standard library.  It verifies the six generator identities
in $\mathbb Q[X]/(X^2+X-1)$, the two primitive own-base relations
at $h=3,4$, all eighteen primitive prime-ideal witnesses, and the
rational restriction ranks for $1\le h\le25$.
The output is preserved in
\texttt{data/golden\_seed\_certificate.json}.
Every calculation in this replay uses exact integer or rational
arithmetic.  The infinite bounds rest on Flatters' theorem and the
prime-ideal argument above; the proof for all $h\ge5$ rests on the
explicit support constraints.

Several further questions are now more sharply delimited.
First, the unresolved golden formulas at weights five through nine
can be approached in the complete eight-argument symbol alphabet;
the missing step is an exact functional-equation or regulator
certificate.  Second, the same uniform $12/24$ cutoff permits complete
seed computations for other real quadratic units, retaining the unit
index congruence when the base is not fundamental.
Third, an effective prime-ideal cutoff for the cubic plastic and
supergolden fields would give an analogous completeness theorem there.
Finally, formalization can begin with the finite arithmetic interface:
quadratic algebra, exact residue orders, integral elimination, and
the tensor/exterior linear-algebra argument.  The cited
primitive-divisor theorem is an explicit separate dependency.
